% -*- coding: utf-8 -*-
% !TEX program = xelatex
\documentclass[plain,noanswer,medmath]{../../template/jnuexam}
\usepackage{../../template/kysx}

\begin{document}


\examtitle{1994}{3}


\loadproblems{1}{1994P1}%载入试卷一


\makepart{填空题}{本题共5小题，每小题3分，满分15分}


\begin{problem}
若$f(x) = \begin{cases}
  \frac{\sin 2x + \e^{2ax} - 1}{x}, & x \ne 0, \\
                                 a, & x = 0
\end{cases}$在$(-\infty , +\infty)$上连续，则$a =$ \fillin{}．
\end{problem}


\begin{problem}
设函数$y = y(x)$由参数方程$\begin{cases}
  x = t - \ln(1 + t), \\
  y = t^3 + t^2
\end{cases}$所确定，则$\frac{\d^2 y}{\dx^2} =$ \fillin{}．
\end{problem}


\begin{problem}
$\frac{\d}{\dx}\left[\int_0^{\cos 3x} f(t)\dt \right] =$ \fillin{}．
\end{problem}


\begin{problem}
$\int x^3\e^{x^2} \dx =$ \fillin{}．
\end{problem}


\begin{problem}
微分方程$y\dx + (x^2 - 4x)\dy = 0$的通解为 \fillin{}．
\end{problem}


\makepart{选择题}{本题共5小题，每小题3分，满分15分}


\begin{problem}
设$\lim_{x \to 0} \frac{\ln(1 + x) - (ax + bx^2)}{x^2} = 2$，则\pickout{}
\begin{abcd}
\item $a = 1,b = - \frac{5}{2}$．		
\item $a = 0,b = - 2$．
\item $a = 0,b = - \frac{5}{2}$．		
\item $a = 1,b = - 2$．
\end{abcd}
\end{problem}


\begin{problem}
设$f(x) = \begin{cases}
  \frac{2}{3}x^3, & x \le 1, \\
             x^2, & x > 1,
\end{cases}$ 则$f(x)$在$x = 1$处\pickout{}
\begin{abcd}
\item 左、右导数都存在．				
\item 左导数存在，但右导数不存在．
\item 左导数不存在，但右导数存在．	
\item 左、右导数都不存在．
\end{abcd}
\end{problem}


\begin{problem}
设$y = f(x)$是满足微分方程$y'' + y' - \e^{\sin x} = 0$的解，且$f'(x_0) = 0$，则$f(x)$在\pickout{}
\begin{abcd}
\item $x_0$的某个邻域内单调增加．		
\item $x_0$的某个邻域内单调减少．
\item $x_0$处取得极小值．					
\item $x_0$处取得极大值．
\end{abcd}
\end{problem}


\begin{problem}
曲线$y = \e^{\frac{1}{x^2}}\arctan \frac{x^2 + x + 1}{(x + 1)(x - 2)}$的渐近线有\pickout{}
\begin{abcd}
\item 1条．        
\item 2条．        
\item 3条．        
\item 4条．
\end{abcd}
\end{problem}


\useproblem{1}{2}{1}%【同试卷一第二(1)题】


\makepart{解答题}{本题共5小题，每小题5分，满分25分}


\begin{problem}
设$y = f(x + y)$，其中$f$具有二阶导数，且其一阶导数不等于$1$，求$\frac{\d^2 y}{\dx^2}$．
\end{problem}


\begin{problem}
计算$\int_0^1x(1 - x^4)^{\frac{3}{2}} \dx$．
\end{problem}


\begin{problem}
计算$\lim_{n \to \infty} \tan^n\left(\frac{\pi}{4} + \frac{2}{n} \right)$．
\end{problem}


\useproblem{1}{3}{3}%【同试卷一第三(3)题】


\begin{problem}
如图，设曲线方程为$y = x^2 + \frac{1}{2}$，梯形$OABC$的面积为$D$，
曲边梯形$OABC$的面积为$D_1$，点$A$的坐标为$(a,0)$（$a > 0$）．证明：$\frac{D}{D_1} < \frac{3}{2}$．
\par\centerline{\includegraphics[width=0.5\linewidth]{1994C1}}
\end{problem}


\makepart{}{本题满分9分}

\begin{problem}
设当$x > 0$时，方程$kx + \frac{1}{x^2} = 1$有且仅有一个解，求$k$的取值范围．
\end{problem}


\makepart{}{本题满分9分}

\begin{problem}
设$y = \frac{x^3 + 4}{x^2}$，
\begin{enumerate}
\item 求函数的增减区间及极值；
\item 求函数图像的凹凸区间及拐点；
\item 求其渐近线；
\item 作出其图形．
\end{enumerate}
\end{problem}


\makepart{}{本题满分9分}

\begin{problem}
求微分方程$y'' + a^2 y = \sin x$的通解，其中常数$a > 0.$
\end{problem}


\makepart{}{本题满分9分}

\begin{problem}
设$f(x)$在$[0,1]$上连续且递减，证明：当$0 < \lambda < 1$时，$\int_0^{\lambda}f(x)\dx \ge \lambda \int_0^1f(x)\dx$．
\end{problem}


\makepart{}{本题满分8分}

\begin{problem}
求曲线$y = 3 - |x^2 - 1|$与$x$轴围成的封闭图形绕直线$y = 3$旋转所得的旋转体体积．
\end{problem}


\end{document}
